\subsection{The groups $C_2\times C_{42}$ and $C_{84}$}
\label{subsec:decic242}

Put $X=X_1(2,42)$. Its genus is $61$, and its rational Jacobian
has rank zero by \cite[Theorem~3.1]{DEHMZ}.
Its projective index is $1152$, so
\cite[(2.1)]{NV} gives
\[
 \operatorname{gon}_{\overline{\mathbb Q}}X\geq
 \frac{325\cdot1152}{32768}>11.
\]
The curve has good reduction at $5$.
Throughout the argument, rational effective divisors may be reducible.

\begin{lemma}\label{lem:decic242-quotient5}
Put $X=X_1(2,42)$, and let $G_9$ be generated by
\[
 a(P,Q)=(21Q,P+Q),\qquad b(P,Q)=(P,25Q).
\]
Then $G_9\simeq C_3\times C_3$ is defined over $\mathbb Q$.
The quotient $Y=X/G_9$ has genus $5$ and the following function
field. Put
\[
 t=u^2,\quad D=t^2+13t+49,\quad N=t^2+5t+1,\quad
 M=t^4+14t^3+63t^2+70t-7.
\]
Adjoin the coordinates $(x_3,y_3)$ of a nonzero point of order $3$
on
\[
 E_u:\quad y^2=x^3-\frac{N}{48D}x+\frac{M}{864D^2}.
\]
Thus $x_3$ satisfies the $3$-division polynomial, and
$y_3^2=x_3^3-Nx_3/(48D)+M/(864D^2)$.
The Jacobian $J_Y(\mathbb Q)$ has rank zero.
\end{lemma}
\begin{proof}
The two displayed transformations commute and have order $3$.
The first cyclically permutes the three nonzero points of order $2$.
The second fixes the full $2$-structure and multiplies the point
of order $7$ by $4$. They generate a group of order $9$.

We give the descent calculation. The Tate normal form for a point
of order $7$ has
\[
 b_7=v^2(v-1),\qquad c_7=v(v-1).
\]
Put $H=v^2-v+1$, $s=(v^3-3v+1)/(v(v-1))$, and $t=s-8$.
The diamond operator of order $3$ sends $v$ to $(v-1)/v$.
Completing the square and depressing the cubic, followed by
scaling by $H$, gives $E_u$ after adjoining $u^2=t$.
The exact invariant identities are
\[
 D=\frac{H^3}{v^2(v-1)^2},\quad
 \frac{c_4}{H^4}=\frac ND,\quad
 \frac{c_6}{H^6}=-\frac M{D^2},\quad
 \frac{\Delta}{H^{12}}=\frac{t}{D^4}.
\]
The square root $u$ records the orientation of the full
$2$-structure, and is fixed by its alternating subgroup.
The point $14Q$ supplies the nonzero point of order $3$.
These functions are invariant under $G_9$.

The resulting field has degree $8$ over $\mathbb Q(u)$.
The two successive equations have degrees $4$ and $2$ and remain
irreducible over $\mathbb F_5(u)$.
Moreover
\[
 j=\frac{N^3D}{t}
\]
has degree $16$ as a function of $u$. Hence the field has degree
$128$ over $\mathbb Q(j)$. The projective index for $X$ is $1152$,
so this is the full invariant field of $G_9$.
The quotient has good reduction at $5$, since $X$ has good
reduction there and $\#G_9=9$ is prime to $5$.
Its normalized finite function field has genus $5$.
Finally $J_Y$ is a quotient of $J_X$, whose rational rank is zero
by \cite[Theorem~3.1]{DEHMZ}.
\end{proof}

\begin{lemma}\label{lem:decic242-quotient5-hecke}
Normalize $J_Y$ at either rational cusp. The reduction image
of $J_Y(\mathbb Q)$ at $5$ is contained in a subgroup
\[
 K\simeq C_2\times C_2\times C_2\times C_{24}
 \subset J_Y(\mathbb F_5),
\]
of index $18$.
\end{lemma}
\begin{proof}
The curve $Y$ has eight geometric cusps. Their fields are
\[
 \mathbb Q,\quad\mathbb Q,\quad\mathbb Q(\sqrt{21}),\quad
 \mathbb Q(\sqrt{-7}),\quad\mathbb Q(\sqrt{-3}),
\]
where each quadratic field contributes two geometric cusps.
To distinguish the branches, put
\[
 \rho=\frac{M}{12ND},\qquad h=\frac{y_3}{x_3-\rho}.
\]
At $u=0$, the two branch pairs have residues
$h^2=1/84$ and $h^2=-1/28$.
At $u=\infty$, they have residues
$h^2=1/4$ and $h^2=-1/12$.
The first three quadratic fields and the two rational cusps follow.
At $5$ the cusp signature is $1^4 2^2$.

The proven finite Jacobian is
\[
 J_Y(\mathbb F_5)=[2,6,6,48].
\]
The finite cusps generate its entire $2$-primary subgroup,
with invariants $[2,2,2,16]$.
The global rational cuspidal divisors have image $[2,2,8]$.
We do not assume that this is the entire rational image.

On $Y$, the diamond operator $\langle11\rangle$ sends
$(u,x_3,y_3)$ to $(u,x_3,-y_3)$.
The operator
\[
 A_{11}=T_{11}-1-11\langle11\rangle
\]
annihilates $J_Y(\mathbb Q)$.
Indeed the Jacobian has rank zero. Its rational torsion injects
at $5$, and has no $11$-primary part since
$\#J_Y(\mathbb F_5)=3456$. Reduction at $11$ is therefore
injective on this torsion, and the Eichler--Shimura relation gives
the assertion.

We compute $A_{11}$ on the finite cusps using their exact global
Frobenius and diamond permutations. Three ordinary closed cubic
places, together with the finite cusps, generate
$J_Y(\mathbb F_5)$. For each cubic place we lift the level
structure over $\mathbb F_{5^9}$ and compute the twelve cyclic
$11$-isogenies. Equal projected closed places are combined before
dividing the norm divisor by $3$; every resulting multiplicity
is checked to be integral.
The images respect every relation of the full finite group.
The exact kernel is $[2,2,2,24]$, and the image is $[3,6]$.
Independent Sage calculations reconstruct all $36$ checked
kernel-polynomial isogenies and verify that their norms exhaust
the respective $11$-division polynomials. They also reconstruct
the full operator and its kernel by integer linear algebra.
Thus the rational image is contained in $K=\ker A_{11}$.
\end{proof}



\begin{lemma}\label{lem:decic242-first-finite-sieves}
If a rational effective divisor of degree $10$ on $X$ has entirely
noncuspidal support and a noncuspidal component in its reduction
at $5$, then, after a rational level automorphism, its noncuspidal
reduction belongs to the $111$ support patterns below.
\end{lemma}
\begin{proof}
The exact noncuspidal closed-point counts in degrees one through
ten over $\mathbb F_5$ are
\[
 (0,0,0,126,648,2080,11376,48816,211776,983196).
\]
The cusp signature is $1^{18},2^9,3^6,6^3$.
The rational level group has order $36$.
Its noncuspidal orbit counts in degrees four through ten are
\[
 (5,18,64,316,1372,5888,27374).
\]
We normalize one summand and retain all actual places for the
other. There are $35037$ single-support patterns. The pair counts are
\[
\begin{array}{c|rrrr}
\text{type}&(4,4)&(4,5)&(4,6)&(5,5)\\\hline
\text{number}&630&3240&10400&11664.
\end{array}
\]
Repetitions are included, and these $60971$ patterns cover all
noncuspidal parts of degree at most $10$.

First use the two maps to $X_1(2,14)$, retaining $(P,3Q)$ or
retaining their images after quotienting by $\langle14Q\rangle$.
The target has rational Jacobian $[2,2,6,18]$, generated by global
rational cuspidal divisors by
\cite[Proposition~4.16 and Table~3]{DEHMZ}.
Its proven finite Jacobian at $5$ is $[2,2,12,36]$.
Every normalized finite cusp class lies in the rational image,
of index $4$.
The two complete normed image tests leave $4647$ patterns.

Next use the genus-$5$ quotient of
Lemma~\ref{lem:decic242-quotient5} and the subgroup
$K=[2,2,2,24]$ of Lemma~\ref{lem:decic242-quotient5-hecke}.
We allow every finite target cusp complement of the required
degree, through degree $6$.
Points with $j=0$ are retained when the chosen chart is not used.
These tests leave $752$ patterns.

Finally use four maps to $X_1(21)$: retain $2Q$, either on $E$
or on each quotient by one of its three nonzero $2$-torsion points.
The rational image is $[364]$ by
\cite[Theorems~3.1 and 4.13, and Table~2]{DEHMZ}.
Its index in the proven finite Jacobian $[2184]$ is $6$.
All finite cusp complements of degrees at most $6$ are included.
The four tests leave $111$ patterns.
Every map includes the residue norm multiplicity $d/e$ when a
source place of degree $d$ maps to one of degree $e$.

The finite-field enumeration, full level orbits and all class
sums are independently checked.
A second native calculation verifies the exact finite and infinite
ideal identities for every supported projected place.
Unsupported CM chart rows are retained throughout.
\end{proof}

% Standalone fragment for the degree-ten X_1(2,42) argument.
% This proof uses no finite-field enumeration or computational certificate.
\begin{lemma}\label{lem:decic242-no-j0}
Let $K$ be a number field, and let $E/K$ be an elliptic curve
with $j(E)=0$. If $E(K)$ contains a subgroup isomorphic to
$C_2\times C_{42}$, then $[K:\mathbb Q]\geq12$.
\end{lemma}
\begin{proof}
Suppose that $[K:\mathbb Q]<12$ and that such a subgroup exists.
Choose $Q\in E(K)$ of order $42$. Write $E$ as $y^2=x^3+B$, where $B\in K^*$.
Full rational $2$-torsion makes the discriminant $-27B^2$ of
$x^3+B$ a square in $K$. Hence
\[
 M=\mathbb Q(\sqrt{-3})\subseteq K,
 \qquad [K:M]\leq5.
\]
Put $\zeta=(-1+\sqrt{-3})/2$. The automorphism
$(x,y)\mapsto(\zeta x,y)$ is defined over $K$, and all
endomorphisms in $\mathbb Z[\zeta]$ are therefore defined over $K$.

We first show that $E$ has good reduction at every prime of $K$
above $7$. Multiplicative reduction is impossible, since $j(E)=0$
is integral. If the reduction were additive, prime-to-$7$ torsion
would inject into the geometric component group, whose order is
at most $4$. This contradicts the point $7Q$ of order $6$.
Thus the reduction is good.

In $\mathbb Z[\zeta]$ we have
\[
 7=\lambda\overline\lambda,
 \qquad \lambda=3+\zeta,
 \qquad \overline\lambda=3+\zeta^2,
 \qquad \lambda+\overline\lambda=5.
\]
The ideals $(\lambda)$ and $(\overline\lambda)$ are distinct
primes of $M$ above $7$. Put $S=6Q$, which has order $7$.
At least one of $\lambda S$ and $\overline\lambda S$ is nonzero,
since their sum is $5S$. Choose
$\alpha\in\{\lambda,\overline\lambda\}$ such that
\[
 R=\overline\alpha S\ne0.
\]
Then $R\in E[\alpha](K)$ has order $7$.

Both primes of $M$ above $7$ have primes above them in $K$.
Hence we may choose a prime $v$ of $K$ above $(\alpha)$,
whichever nonzero component $R$ was chosen. Put $L=K_v$.
We normalize its valuation by $v_L(L^*)=\mathbb Z$, and write
$e=v_L(7)$. Since $7$ splits in $M$,
\[
 e\leq[L:\mathbb Q_7]\leq[K:M]\leq5.
\]
The isogeny $[\alpha]$ has degree $7$. Its reduction also has
degree $7$, and its differential is zero, since its multiplier
on invariant differentials is $\alpha$. Hence its reduction is
purely inseparable. It follows that $R$ reduces to the identity
and belongs to the formal group of $E/L$.
The conjugate isogeny $[\overline\alpha]$ has unit differential
at $v$. Its reduction is separable of degree $7$, so the good
reduction is ordinary.

Let $t$ be the formal parameter of $R$. Then $t\ne0$ and
$s=v_L(t)$ is a positive integer. For the ordinary formal group,
write
\[
 [7](T)=7T+a_2T^2+\cdots+a_6T^6+a_7T^7+\cdots.
\]
Here $a_i\in7\mathcal O_L$ for $2\leq i\leq6$, and
$a_7\in\mathcal O_L^*$.
Indeed, the identity $[7]^*\omega=7\omega$ for an integral
invariant differential gives $[7]^{\prime}(T)\in7\mathcal O_L[[T]]$.
This proves the divisibility of $a_2,\ldots,a_6$; the unit
condition on $a_7$ is the height-one condition.
In the equality $[7](t)=0$, the only terms which can have minimal
valuation are $7t$ and $a_7t^7$. Consequently
\[
 e+s=7s,\qquad e=6s\geq6.
\]
This contradicts $e\leq5$.
\end{proof}

% Companion to CM_j0_degree10.tex. No odd-degree parity argument is used.
\begin{lemma}\label{lem:decic242-order3-stabilizers}
Put $X=X_1(2,42)$. Let $\overline N$ be a noncuspidal effective
divisor of degree $9$ on $X_{\mathbb F_5}$, fixed by a rational
level automorphism $h$ of order $3$.
There is no effective divisor $D$ on $X$, defined over $\mathbb Q$,
of degree $10$, with entirely noncuspidal support and reduction
$\overline N+\overline c$, where $\overline c$ is a cusp.
\end{lemma}
\begin{proof}
The rational Jacobian of $X$ has rank zero by
\cite[Theorem~3.1]{DEHMZ}.
Its projective index is $1152$, so the bound in
\cite[(2.1)]{NV} gives
\[
 \operatorname{gon}_{\overline{\mathbb Q}}X
 \geq\frac{325\cdot1152}{32768}>11.
\]
Suppose that $D$ exists. Lift $\overline c$ to a cusp $c$ over
an extension unramified at $5$.
The class
\[
 [hD-D-hc+c]
\]
is torsion by rank zero and Manin--Drinfeld. Its reduction is
zero, since $h\overline N=\overline N$.
Injectivity of torsion reduction over an unramified extension of
$\mathbb Q_5$ gives
\[
 hD+c\sim D+hc.
\]
These are effective divisors of degree $11$, below geometric
gonality. Hence they are equal.
Comparing their noncuspidal parts gives $hD=D$.

A geometric fixed point of $h$ would give an elliptic
automorphism of order $3$ or $6$ and would therefore have
$j$-invariant $0$.
Every closed point in the support of $D$ has degree at most $10$.
Lemma~\ref{lem:decic242-no-j0} excludes such a point with
$j$-invariant $0$. Thus $h$ acts freely on the geometric support
of $D$. Its orbits have length $3$, and $hD=D$ makes the
multiplicities constant on each orbit. It follows that $3$
divides $\deg D=10$, a contradiction.
\end{proof}

\noindent\emph{Application to the three degree-$9$ supports.}
Let
\[
 k=\mathbb F_5[\alpha]/(\alpha^9+2\alpha^3+\alpha+3).
\]
For $B=\sum_{i=0}^8B_i\alpha^i$ and
$C=\sum_{i=0}^8C_i\alpha^i$, use the level structure
\[
 E:\ y^2=x^3+Cx^2+B(1-B-C)x,
 \qquad P=(0,0),\quad Q=(B,B).
\]
The three remaining degree-$9$ supports have coefficient vectors
\[
\begin{array}{c|c|c}
\text{case}&(B_i)&(C_i)\\ \hline
1777&(0,0,0,0,2,0,2,1,4)&(4,4,3,1,0,0,4,2,2)\\
1780&(0,0,1,0,1,3,4,1,2)&(3,4,1,0,4,1,1,4,4)\\
1783&(0,0,0,1,0,1,0,2,3)&(4,2,3,2,3,3,0,2,1).
\end{array}
\]
Put
\[
 h_1(P,Q)=(21Q,P+25Q),\qquad
 h_2(P,Q)=(P+21Q,P+4Q).
\]
Both have order $3$. On case $1777$, $h_1$ agrees with the
third power of Frobenius. On cases $1780$ and $1783$, $h_2$
agrees with the third and sixth powers of Frobenius,
respectively. Each closed degree-$9$ divisor is therefore fixed
by the indicated automorphism, and
Lemma~\ref{lem:decic242-order3-stabilizers} excludes all three.

The complete automorphism and cusp-action enumeration is in
\path{next242_cheap_stabilizers.py}; it checks all $36$ level
automorphisms on all $111$ input supports and all $72$ geometric
cusps. It removes exactly these three cases.
The three displayed Frobenius identities and the exact order-$3$
level identities are independently verified in Magma by
\path{next242_order3_phases.m}; the output is
\path{degree9_order3_native_phases.m}.
The remaining $108$ supports, with their original case IDs and
field polynomials, are stored in
\path{after_cheap_stabilizers_unified.m}.


\begin{lemma}\label{lem:decic242-X042-image}
Let $q:X_1(42)\to Z=X_0(42)$ be the natural map.
Normalize at a rational cusp of $X_1(42)$ and its image.
The reduction of $q_*J_1(42)(\mathbb Q)$ at $5$ is a subgroup
$H=[2,48]$ of index $192$ in
\[
 J_Z(\mathbb F_5)=[2,4,48,48].
\]
The classes of finite effective cuspidal divisors of degrees
$0,\ldots,6$, modulo $H$, form sets of cardinalities
\[
 (1,6,12,16,20,24,24).
\]
\end{lemma}
\begin{proof}
A convenient model of $Z$ is the normalization of
\[
 \frac{N(t)^3D(t)}t=j(E_z),\qquad
 E_z:\ y^2+(1-z)xy-(z+z^2)y=x^3-(z+z^2)x^2,
\]
where $D=t^2+13t+49$ and $N=t^2+5t+1$.
The equation has degree $8$ in $t$, is irreducible, and has genus
$5$. It is the fibre product of $X_0(7)$ and
$X_1(6)=X_0(6)$ over the $j$-line.

We use global rational cuspidal divisors on $X_1(42)$.
These generate its rational Jacobian by
\cite[Theorems~3.1 and 4.13, and Table~2]{DEHMZ}.
Here the rational cuspidal group is defined as the image of
Galois-invariant cuspidal \emph{divisors}, not merely the
Galois-fixed part of the geometric cuspidal class group.
In particular, the stated theorem supplies the required integral
Galois descent.

All eight cusps of $Z$ are rational.
A cusp with denominator $g\mid42$ receives the global cusp
orbits on $X_1(42)$ in the following table:
\[
\begin{array}{c|rrrrrrrr}
g&1&2&3&6&7&14&21&42\\\hline
\text{number of global orbits}&6&6&3&3&1&1&1&1\\
\text{degree of each orbit}&1&1&2&2&6&6&6&6.
\end{array}
\]
Thus the pushforward generators, normalized at the chosen base,
are the corresponding cusp classes multiplied by the last row.
Their exact subgroup in the proven finite Jacobian is $[2,48]$.
No assertion that this is the whole of $J_Z(\mathbb Q)$ is needed.

We check this pushforward directly from the universal level
structure on $X_1(42)$.
For the point $Q$ of order $42$, let $z$ be the Tate parameter
for $7Q$ of order $6$.
For $6Q$ of order $7$, write
$b_7=v^2(v-1)$ and $c_7=v(v-1)$; then
\[
 t=\frac{v^3-3v+1}{v(v-1)}-8.
\]
Valuations of $z,z+1,9z+1,t$ distinguish the eight target cusps.
All $22$ global source orbits have the stated pushforwards, and
each target cusp has six geometric preimages.
An independent Tate-label calculation reproduces every global
and finite cusp orbit and these valuation signatures.
Finally we enumerate all effective target cusp complements
through degree $6$. This gives the displayed cardinalities.
\end{proof}

\noindent
To evaluate the normalized model at all noncuspidal places, we
also use the rational function
\[
 \eta=-\frac{c_6(E_z)N(t)D(t)}{c_4(E_z)M(t)},\qquad
 M(t)=t^4+14t^3+63t^2+70t-7.
\]
If $k_6,k_7$ are the scales in the two Tate normal forms, then
$\eta=(k_7(v^2-v+1)/k_6)^2$.
This expression evaluates the function without discarding the
branches above $j=0$ or $1728$.
The three coordinates $(z,t,\eta)$ determine a unique normalized
place in every finite test below.


\begin{lemma}\label{lem:decic242-X042-refinement}
The $108$ supports left by
Lemma~\ref{lem:decic242-order3-stabilizers} reduce to ten after
the necessary image tests on $X_0(42)$.
\end{lemma}
\begin{proof}
For the original curve $E$, retain $Q$, $P+Q$ or $P+22Q$.
The other three maps first quotient by
$\langle P\rangle$, $\langle P+21Q\rangle$ or
$\langle21Q\rangle$.
For the first two kernels retain the image of $Q$; for the last
retain $\phi(P)+\phi(Q)$.
All six points have exact order $42$, and hence define rational
maps to $X_1(42)$ and then to $X_0(42)$.

Write the reduction as $N+C$ and let $f$ be one of the six maps.
For the quotient $G/H$ in Lemma~\ref{lem:decic242-X042-image},
let $\mathcal C_e$ be the set of normalized finite cusp classes
of degree $e$. A necessary condition is
\[
 -[f_*N-(\deg N)c_0]\in\mathcal C_{10-\deg N}\pmod H.
\]
The six tests leave ten patterns: one of type $(6)$, one of type
$(10)$, four of type $(4,4)$, two of type $(4,6)$ and two of type
$(5,5)$.
The normalized coordinates above are used at every CM residue
branch; no point is omitted because of a chart pole.
Independent arithmetic verifies all $211$ projections, including
$45$ explicit $2$-isogenies, the global cusp norms and all class
conditions. Every class calculation has a second exact native
ideal-identity replay.
\end{proof}

\begin{lemma}\label{lem:decic242-hecke17}
Put $Z=X_1(42)$ and
$A_{17}=T_{17}-1-17\langle17\rangle$.
The operator $A_{17}$ annihilates every cuspidal divisor on
$Z_{\mathbb F_5}$ defined over $\mathbb F_5$.
If a rational effective divisor on $Z$ reduces to $N+C$, where
$N$ is noncuspidal and $C$ is cuspidal, then $A_{17}N$ is
principal.
\end{lemma}
\begin{proof}
Write a geometric cusp in Tate coordinates as $(a,b)$, where
the point of order $42$ is represented by
$\zeta_{42}^a q^{b/42}$.
Galois and diamond operators act by
\[
 \sigma_u(a,b)=(ua,b),\qquad
 \langle u\rangle(a,b)=(ua,ub).
\]
For a prime $\ell\nmid42$, the quotient by the toric kernel
$\mu_\ell$ gives the first term in the boundary formula
\[
 T_\ell c=\sigma_\ell c+
 \ell\langle\ell\rangle\sigma_\ell^{-1}c.
\]
The other $\ell$ kernels give the second term.
These conventions agree with retaining the image of the marked
point under each isogeny.

At $5$, Galois on the cusps is generated by $\sigma_5$.
Since $17\equiv5^5\pmod{42}$, every $\mathbb F_5$-rational
cuspidal divisor $C$ satisfies
$\sigma_{17}C=C=\sigma_{17}^{-1}C$.
The boundary formula gives $A_{17}C=0$ as a divisor.
This applies to arbitrary multiplicities and does not require
$C$ to lift to a rational cuspidal divisor over $\mathbb Q$.

The group $J_1(42)(\mathbb Q)$ is finite by
\cite[Theorems~3.1 and~4.13, Table~2]{DEHMZ}.
The Eichler--Shimura relation at $17$, together with injectivity
of torsion reduction at this odd good prime, gives
\[
 A_{17}J_1(42)(\mathbb Q)=0.
\]
Choose a rational cusp $c_0$ on $Z$. The boundary formula gives
$A_{17}c_0=0$ as a divisor.
Consequently $A_{17}D$ is principal for every rational divisor
$D$ on $Z$.
Specializing at $5$ gives that $A_{17}(N+C)$ is principal.
Since $A_{17}C=0$, the divisor $A_{17}N$ is principal.
\end{proof}

\noindent\emph{Exact finite-field test.}
For each source support and each chosen map to $Z$, we enumerate
every Frobenius orbit of cyclic kernels of order $17$.
The orbit lengths sum to $18$, and the corresponding factors
exhaust the degree-$144$ $17$-division polynomial in the
$x$-coordinate. Every quotient point is multiplied by its exact
residue-degree ratio. Thus the positive Hecke divisor has degree
$18\deg N$.
The resulting divisor
\[
 T_{17}N-N-17\langle17\rangle N
\]
has degree zero. A Riemann--Roch dimension of zero contradicts
Lemma~\ref{lem:decic242-hecke17}.
Both maximal-order ideal bases are retained for an independent
Riemann--Roch replay.
For supports without a cusp complement the same test may use
$A_{11}=T_{11}-1-11\langle11\rangle$.


\begin{lemma}\label{lem:decic242-final-norm-sieve}
Nine of these ten patterns do not arise from a rational effective
divisor of degree $10$ with noncuspidal support.
The remaining pattern is case $26$, a degree-$6$ noncuspidal
place and a degree-$4$ cuspidal complement.
\end{lemma}
\begin{proof}
For the five patterns whose noncuspidal part already has degree
$10$, use $A_{11}$ on $X_1(42)$.
For the four patterns of type $(4,4)$, use $A_{17}$ from
Lemma~\ref{lem:decic242-hecke17}.
The first retaining map suffices except for case $35579$, where
we retain $P+Q$.
In each case the resulting degree-zero Hecke divisor has
Riemann--Roch dimension zero, contradicting principality.

The independent reconstruction checks $16$ normed source
summands and $78$ kernel orbits, accounting for $240$ geometric
cyclic kernels of orders $11$ and $17$.
Their norm products exhaust the appropriate division polynomials.
All quotient curves, marked images and residue-degree weights
are verified independently.
A second native calculation compares both exported ideal-module
lattices with the native ideals and reproduces all nine zero
Riemann--Roch dimensions.
The only pattern not excluded is case $26$.
\end{proof}

\begin{lemma}\label{lem:decic242-case26-cubic}
In case $26$, the degree-$4$ cusp complement contains one closed
cubic cusp and one rational cusp.
\end{lemma}
\begin{proof}
Use the first map to $X_0(42)$ in the exact quotient sieve.
Its normalized noncuspidal class is
\[
 v=(1,3,6,12,0)\in
 G=[2,4,48,48,0].
\]
The rational-image subgroup $H=[2,48]$ has generators whose
first coordinates are all even. The first coordinate modulo $2$
therefore vanishes on $H$, whereas it takes the value $1$ on $v$.
The residue cusp complement cannot be the reduction of a rational
cuspidal divisor: its pushforward would contribute a class in $H$
and could not correct $v$.

The global source cusp signature is $1^{18}2^9 6^6$, and the
signature at $5$ is $1^{18}2^9 3^6 6^3$.
Every rational or quadratic residue cusp is globally intact.
Hence every degree-$4$ cusp complement without a cubic cusp lifts
to a rational cuspidal divisor. This is impossible by the preceding
class test. The only remaining pattern is one cubic cusp and one
rational cusp.
\end{proof}

The exact subgroup vectors and noncuspidal class are in
\path{geometry/x0_42_proven_picard_v2.m} and
\path{geometry/x0_42_full_arithmetic.m}.
The global cusp pushforwards are certified in
\path{geometry/x0_42_direct_cusp_pushforward.m}.
An independent integer enumeration verifies both source cusp
signatures and the intact rational and quadratic cusps; see
\path{case26_cubic_complement.json}.

\begin{lemma}\label{lem:decic242-last-support}
The remaining support in case $26$ cannot be the reduction of a
rational effective divisor of degree $10$ with entirely
noncuspidal support.
\end{lemma}
\begin{proof}
Write its noncuspidal part as $\overline N$.
For an exact identification, let
\[
 k=\mathbb F_5[\alpha]/
 (\alpha^6+\alpha^4+4\alpha^3+\alpha^2+2).
\]
The coefficient vectors of $B,C\in k$, in ascending powers of
$\alpha$, are
\[
 (B_i)=(0,0,2,4,2,3),\qquad(C_i)=(4,1,4,0,0,1).
\]
The level structure
\[
 E:\ y^2=x^3+Cx^2+B(1-B-C)x,\qquad P=(0,0),\quad Q=(B,B)
\]
represents $\overline N$.
Put
\[
 h(P,Q)=(P+21Q,P+4Q).
\]
This automorphism has order $3$ and agrees with the second power
of Frobenius on $\overline N$.
It acts freely on all $72$ geometric cusps.

Let $\pi:X\to W=X/\langle h\rangle$.
The quotient has genus $17$, and
$\overline N=\pi^*q$ for a closed quadratic place $q$.
We describe the model and the section calculation used below.
With the notation of Lemma~\ref{lem:decic242-quotient5}, adjoining
\[
 w^3-\frac{N(t)}{48}w+
 \frac{M(t)(2t+13)}{1728D(t)}-\frac{u}{8D(t)}=0
\]
to the genus-$5$ field gives the field of $W$.
Indeed, if $r_0,r_1,r_2$ are the ordered normalized $2$-division
roots and $v_0=v$, $v_1=(v-1)/v$, $v_2=1/(1-v)$, then
$w=r_0v_0+r_1v_1+r_2v_2$ satisfies this equation.
The generator used for this invariant-field model is
$(21Q,P+37Q)=h^2$; it agrees with the fourth power of Frobenius
on $\overline N$. Thus it gives the same quotient.

Over $\mathbb F_5$, the genus-zero $(u,w)$-field has a parameter
$z$ with
\[
 u=\frac{z^3+3z^2+4z+3}{z^2+1}.
\]
Adjoining the coordinates of the nonzero $3$-point on $E_u$
gives an absolute degree-$8$ model over $\mathbb F_5(z)$ of
genus $17$.
The source function $u$ is regular on $\overline N$ and satisfies
$u^2+4u+2=0$ there.
We retain \emph{all} quadratic places of $W$ above this polynomial.
There are two. Both are unramified with inert cubic fibre, so
they form a containing set for the actual place $q$.
This avoids evaluating a chart with poles at the CM residue point.

The source cusp signature is $1^{18},2^9,3^6,6^3$.
By Lemma~\ref{lem:decic242-case26-cubic}, the complement
$\overline C$ consists of one cubic and one rational cusp.
The quotient cusp signature is $1^6,2^3,3^2,6$.
The $108$ possible source complements give exactly $12$ divisors
$\pi_*\overline C$, one for each choice of a cubic and a rational
cusp of $W$; each choice has nine source preimages.

We compute the required section dimensions after extending
constants to $\mathbb F_{25}$.
Put $s=u^2+8$, $D=s^2-3s+9$, and
$A=2s^3-9s^2+27s-27$.
Adjoining the Tate parameter $v$ recovers the cyclic degree-$3$
source cover:
\[
 v^3-sv^2+(s-3)v+1=0.
\]
For a primitive cube root $\zeta\in\mathbb F_{25}$, this is the
Kummer cover $y^3=f$, where
\[
 f=\frac{A+3(\zeta-\zeta^2)D}{2}.
\]
The identities
$A^2-4D^3=-27D^2$ and $3v-s=y+D/y$ verify the correspondence.
For $j=0,1,2$, put
\[
 D_j=\left\lfloor\frac{j\operatorname{div}(f)}3\right\rfloor.
\]
Their degrees are $0,-6,-6$, and the branch divisor has degree
$12$. The usual integral local bases of a tame Kummer cover give
\[
 h^0_X(\pi^*T)=\sum_{j=0}^2h^0_W(T+D_j).
\]

A quadratic place over $\mathbb F_5$ becomes a sum of two rational
places over $\mathbb F_{25}$.
The degree-$6$ source place becomes a sum of two cubic places;
we retain both summands in each divisor.
For each of the two admitted $q$ and all $12$ cusp choices, take
$T=q+\pi_*\overline C$, of degree $6$.
The $24$ exact calculations, comprising $72$ Riemann--Roch spaces,
give
\[
 (h^0(T+D_0),h^0(T+D_1),h^0(T+D_2))=(1,0,0).
\]
Consequently
\begin{equation}\label{eq:decic242-final-section}
 h^0_X(\overline N+\pi^*\pi_*\overline C)=1
\end{equation}
for every compatible complement.
All places are exported by two exact ideal generators.
The field is an absolute extension of degree $8$ over the rational
function field; its finite and infinite ideal module bases each
have eight elements. A second native calculation reconstructs
the character divisors, compares both module lattices in both
directions, and reproduces all $72$ dimensions.
Independent arithmetic also verifies the cyclic-cover identity,
the Frobenius phases and the complete cusp-orbit coverage.

Suppose that a rational noncuspidal divisor $D$ of degree $10$
has the proposed reduction.
Lift $\overline C$ to a cuspidal divisor $C$ over an unramified
local field at $5$.
The class $[hD-D-hC+C]$ is torsion by rank zero and
Manin--Drinfeld, and reduces to zero.
Odd unramified torsion injectivity makes it zero.
Thus $L=[D-C]$ is $h$-invariant.
Put
\[
 \Omega=C+hC+h^2C=\pi^*\pi_*C.
\]
The effective divisors $D+\Omega-C$ and $hD+\Omega-hC$ are
linearly equivalent.
By \eqref{eq:decic242-final-section} and upper semicontinuity,
their common line bundle has at most one section, so the divisors
are equal. Comparing their cuspidal parts gives $C=hC$.
Since $h$ acts freely on the geometric cusps, this forces
$3\mid\deg C=4$, a contradiction.
\end{proof}


\begin{proposition}\label{prop:decic242-excluded}
No elliptic curve over a number field of degree $10$ has a
subgroup isomorphic to $C_2\times C_{42}$ or $C_{84}$.
\end{proposition}
\begin{proof}
A hypothetical level structure gives a rational effective
noncuspidal divisor $D$ of degree $10$ on $X$.
If its reduction is entirely cuspidal, lift that reduction to a
cuspidal divisor $C$ over an unramified extension at $5$.
Rank zero and Manin--Drinfeld make $[D-C]$ torsion.
Its reduction is zero, so odd unramified torsion injectivity makes
it principal. The two distinct effective divisors have degree
$10$, below geometric gonality, a contradiction.

Otherwise the finite sieves, the three order-$3$ stabilizer
exclusions and the final norm tests leave only case $26$.
Lemma~\ref{lem:decic242-last-support} excludes this case.
Thus $C_2\times C_{42}$ is impossible.

Finally a point $R$ of order $84$ gives such a subgroup after
quotienting by $\langle42R\rangle$.
The image $Q$ has order $42$.
The nonzero point $S$ in the dual kernel is rational and has
order $2$. It differs from $21Q$, since the dual isogeny kills
$S$ but sends $21Q$ to $42R\ne0$.
This excludes $C_{84}$.
\end{proof}
